\chapter{Twierdzenie Sarda i transwersalność}
\label{ch:sard-transwersalnosc}
\index{transwersalnosc@Transwersalność}

W Rozdziale~\ref{ch:rozmaitosci} poznaliśmy przeciwobrazy
wartości regularnych, a w Rozdziale~\ref{ch:geodezyjne-exp}
geodezyjne i ich wariacje. Tutaj odpowiemy na dwa pytania:
\emph{czy można znaleźć wartości regularne} i
\emph{czy dwa obiekty można ustawić w położeniu ogólnym?}
Twierdzenie Sarda odpowiada na pierwsze, a transwersalność
precyzuje drugie. Będą one potrzebne do późniejszego
badania punktów krytycznych, przecięć rozmaitości
stabilnych i niestabilnych, włożeń oraz zabiegów na
rozmaitościach. Nazwy z przyszłej teorii Morse'a i
Smale'a pojawią się na końcu jako zastosowanie aparatu;
ich właściwa teoria będzie osobnym rozdziałem.

O ile nie zaznaczymy inaczej, rozmaitości i podrozmaitości
są gładkie i bez brzegu, Hausdorffa i drugoprzeliczalne.
Podrozmaitość rozumiemy, zgodnie z Rozdziałem~\ref{ch:rozmaitosci},
z topologią podprzestrzeni i mapami plasterkowymi
(ang.\ \emph{embedded submanifold}); nie zakładamy jej
domkniętości, chyba że jest to wyraźnie powiedziane.
Przypadek dziedziny z brzegiem omówimy osobno.

\section{Wartości krytyczne i twierdzenie Sarda}
\label{sec:sard}

Przypomnijmy Definicję wartości regularnej z
Rozdziału~\ref{ch:rozmaitosci}: $y\in N$ jest regularna
dla $f:M\to N$, gdy $\dd f_x:T_xM\to T_yN$ jest
surjekcją dla \emph{każdego} $x\in f^{-1}(y)$.
Jeżeli $f^{-1}(y)=\varnothing$, warunek jest spełniony.
Przeciwna sytuacja wymaga rozróżnienia punktu dziedziny
od jego obrazu.

\begin{definicja}[Punkt i wartość krytyczna]
\index{punkt i wartosc krytyczna@Punkt i wartość krytyczna}
\label{def:sard-krytyczne}
Niech $f:M^m\to N^n$ będzie gładkie.
Punkt $x\in M$ jest \emph{krytyczny}, gdy
$\operatorname{rank}\dd f_x<n$; jego obraz $f(x)$
nazywamy \emph{wartością krytyczną}. Zbiór punktów
krytycznych oznaczamy przez $C(f)$, a zbiór wartości
krytycznych przez $f(C(f))$. Jeśli $n=1$, punkt
krytyczny funkcji $f:M\to\R$ spełnia $\dd f_x=0$.
\end{definicja}

Przykładowo $f(x,y)=x^2+y^2$ ma jedyny punkt
krytyczny $(0,0)$ i jedyną wartość krytyczną $0$.
Poziomice $f^{-1}(c)$ dla $c>0$ są okręgami, poziomica
$c=0$ jest pojedynczym punktem, a dla $c<0$
przeciwobraz jest pusty. Pusty przeciwobraz też
odpowiada wartości regularnej; nie należy utożsamiać
„regularna” z „poziomica niepusta”.

\begin{figure}[htbp]
\centering
\begin{tikzpicture}[>=Latex,font=\small,line cap=round]
 \begin{scope}[shift={(-2.5,0)}]
  \draw[->,manifoldblue!60] (-1.8,0)--(1.85,0) node[right] {$x$};
  \draw[->,manifoldblue!60] (0,-1.7)--(0,1.8) node[above] {$y$};
  \draw[manifoldblue!75!black,thick] (0,0) circle (.73);
  \draw[manifoldgreen!75!black,thick] (0,0) circle (1.38);
  \fill[manifoldred] (0,0) circle (2.3pt);
  \node[manifoldblue!75!black] at (-1.16,.88) {$f=c_1$};
  \node[manifoldgreen!65!black] at (1.40,1.26) {$f=c_2$};
 \end{scope}
 \begin{scope}[shift={(3,0)}]
  \draw[->,manifoldblue!75!black,thick] (0,-1.6)--(0,1.65)
    node[above] {$\R$};
  \fill[manifoldred] (0,-.97) circle (2.3pt)
    node[right=5pt] {$0$: wartość krytyczna};
  \fill[manifoldblue] (0,-.19) circle (2.1pt)
    node[right=5pt] {$c_1$: wartość regularna};
  \fill[manifoldgreen!75!black] (0,.85) circle (2.1pt)
    node[right=5pt] {$c_2$: wartość regularna};
 \end{scope}
\end{tikzpicture}
\caption{Dla $f(x,y)=x^2+y^2$ punkt krytyczny leży w
dziedzinie, a jego wartość krytyczna na osi wyników.
Poziomice dodatnie są gładkimi okręgami. Schemat
rozróżnia oba pojęcia.}
\label{fig:sard-poziomice}
\end{figure}

\begin{definicja}[Zbiór miary zero na rozmaitości]
\index{zbior miary zero na rozmaitosci@Zbiór miary zero na rozmaitości}
\label{def:sard-zero-miary}
Podzbiór $A\subset\R^n$ ma \emph{miarę zero},
jeśli dla każdego $\varepsilon>0$ można pokryć go
przeliczalnie wieloma kostkami, których suma
$n$-wymiarowych objętości jest mniejsza niż
$\varepsilon$. Zbiór $B\subset N^n$ ma miarę zero,
jeśli dla każdej mapy $(V,\psi)$ rozmaitości $N$
zbiór $\psi(B\cap V)$ ma miarę zero w $\R^n$.
\end{definicja}

Ta definicja nie zależy od mapy. Istotnie, przejście
między dwiema mapami jest gładkim dyfeomorfizmem.
Na każdej zwartej kostce zawartej w dziedzinie przejścia
pochodna jest ograniczona, więc twierdzenie o wartości
średniej daje tam stałą Lipschitza $L$. Jeśli mała kostka
pokrywająca zbiór zerowy ma bok $r$, obraz jej przecięcia
z ustaloną zwartą kostką mieści się w kostce o boku $2L\sqrt n\,r$.
Suma objętości rośnie zatem najwyżej o stały czynnik.
Pokrycia o dowolnie małej sumie objętości pokazują,
że obraz zbioru miary zero znów ma miarę zero. Dziedzinę przejścia można
pokryć przeliczalnie wieloma zwartymi kostkami.
Drugi aksjomat przeliczalności zapewnia także
przeliczalne pokrycie rozmaitości mapami. Z tego
powodu lokalne dowody o mierze zero można
sumować po mapach.

\begin{twierdzenie}[Sarda, wersja gładka]
\label{thm:sard}
Niech $f:M^m\to N^n$ będzie gładkim odwzorowaniem
rozmaitości bez brzegu, $n\geq1$. Zbiór jego
\emph{wartości krytycznych} $f(C(f))$ ma miarę zero
w $N$. W szczególności wartości regularne są gęste
w $N$.
\end{twierdzenie}

\begin{proof}
Wyjaśnimy wszystkie kroki dowodu dla $C^\infty$.
Wybieramy mapy dziedziny tak małe, aby ich obrazy przez
$f$ mieściły się w mapach przeciwdziedziny. Wystarczy więc
rozważyć $f:U\subset\R^m\to\R^n$, z otwartym $U$. Zbiory
$C(f)\cap K$ dla zwartych kostek $K\Subset U$
są zwarte: warunek $\operatorname{rank}\dd f<n$
oznacza znikanie wszystkich minorów rzędu $n$
(jeśli $m<n$, zachodzi automatycznie). Ich obrazy
są więc zwarte, a $f(C(f))$ jest lokalnie
przeliczalną sumą zbiorów zwartych. Wolno
zastosować do nich twierdzenie Fubiniego:
jeżeli każda sekcja takiego mierzalnego zbioru
przez hiperpowierzchnię $\{y_1=t\}$ ma
$(n-1)$-wymiarową miarę zero, to cały zbiór
ma $n$-wymiarową miarę zero.

\emph{Krok 1: źródło mniejszego wymiaru.}
Jeśli $m<n$, na zwartej kostce $Q\Subset U$
odwzorowanie $f$ jest lipschitzowskie.
Podzielmy $Q$ na $k^m$ kostek o boku $a/k$.
Obraz każdej mieści się w kostce o boku
$Ca/k$, zatem suma objętości pokrycia
$f(Q)$ jest co najwyżej
$C'k^m k^{-n}=C'k^{m-n}\to0$.
Stąd cały obraz $f(U)$, a więc tym bardziej
zbiór wartości krytycznych, ma miarę zero.
To daje również początek indukcji po $m$;
dla $m=0$ obraz drugoprzeliczalnej
rozmaitości zerowymiarowej jest przeliczalny.

\emph{Krok 2: punkty krytyczne dodatniego rzędu.}
Załóżmy twierdzenie dla wszystkich odwzorowań
o dziedzinie wymiaru $m-1$ i niech
$x\in C(f)$, ale $\operatorname{rank}\dd f_x>0$.
Po przestawieniu współrzędnych
$\partial f_1/\partial x_1(x)\ne0$.
Twierdzenie o funkcji odwrotnej zastosowane do
\[
 \Phi(x_1,\ldots,x_m)
       =(f_1(x),x_2,\ldots,x_m)
\]
daje dyfeomorfizm pewnego otwartego otoczenia $U_{\mathrm{lok}}$
na jego obraz. W nowych współrzędnych
przepis $f$ ma postać
$g(t,u)=(t,h(t,u))$, gdzie $u\in\R^{m-1}$.
Jeśli $n=1$, ta postać jest submersją,
więc rozważany punkt nie byłby krytyczny.
Przy $n\geq2$ macierz pochodnej ma bloki
\[
 \dd g_{(t,u)}
  =\begin{pmatrix}
       1&0\\ *&\dd_u h(t,u)
    \end{pmatrix},
 \qquad
 \operatorname{rank}\dd g
      =1+\operatorname{rank}\dd h_t(u),
 \quad h_t(u):=h(t,u).
\]
Zatem w sekcji $\{y_1=t\}$ obrazy punktów
krytycznych $g$ pochodzą z wartości krytycznych
$h_t:U_t\to\R^{n-1}$, gdzie
$U_t=\{u:(t,u)\in\Phi(U_{\mathrm{lok}})\}$ jest otwartym
przekrojem nowej dziedziny. Indukcja mówi,
że taka sekcja ma miarę zero. Obraz lokalnego
zbioru punktów krytycznych jest mierzalny
(po wyczerpaniu otwartego otoczenia kostkami
jest sumą obrazów zwartych zbiorów).
Fubini daje miarę zero całego obrazu.
Przeliczalne lokalne otoczenia pokrywają
wszystkie punkty krytyczne dodatniego rzędu.

\emph{Krok 3: punkty o znikających pochodnych.}
Niech $D_k\subset U$ oznacza zbiór punktów,
w których znikają wszystkie pochodne cząstkowe
składowych $f$ rzędów $1,\ldots,k$.
Mamy $D_1\subset C(f)$ oraz
$D_{k+1}\subset D_k$. Jeżeli
$x\in D_k\setminus D_{k+1}$, to pewna pochodna
rzędu $k+1$ jest niezerowa. Wybierzmy
pochodną rzędu $k$ jednej składowej, oznaczoną
przez $w$, której pewna pochodna pierwsza
nie znika w $x$. Ponieważ $w=0$ na $D_k$,
twierdzenie o funkcji uwikłanej mówi, że
$H=\{w=0\}$ jest w otoczeniu $x$ gładką
hiperpowierzchnią zawierającą $D_k$.
Każdy punkt $D_k\cap H$ jest krytyczny dla
$f|_H:H^{m-1}\to\R^n$: już $\dd f=0$
na $D_k$. Na mocy indukcji jego obrazy
tworzą zbiór miary zero. Przeliczalne
pokrycie punktów $D_k\setminus D_{k+1}$
daje tę samą tezę dla całej tej warstwy.

\emph{Krok 4: punkty płaskie do wysokiego rzędu.}
Wybierzmy $k\geq1$ tak duże, że
$n(k+1)>m$. Na zwartej kostce $Q\Subset U$
pochodne rzędu $k+1$ są ograniczone.
Wzór Taylora wokół dowolnego $x\in D_k\cap Q$
ma zerowe składniki rzędów $1,\ldots,k$,
więc dla $y$ w tej samej małej kostce
\[
 |f(y)-f(x)|\leq C|y-x|^{k+1}.
\]
Podzielmy $Q$ na $j^m$ kostek o boku $a/j$.
W każdej, która przecina $D_k$, wybierzmy
takie $x$. Obraz jej punktów należących
do $D_k$ mieści się w kuli promienia
$C'j^{-(k+1)}$. Pokrywając każdą kulę
kostką, dostajemy sumę objętości
co najwyżej $C''j^{m-n(k+1)}\to0$.
Zatem $f(D_k\cap Q)$ ma miarę zero;
liczymy po kostkach wyczerpujących $U$.

Na koniec rozkładamy zbiór krytyczny na
$C(f)\setminus D_1$, warstwy
$D_i\setminus D_{i+1}$ dla $1\leq i<k$
oraz $D_k$. Ich obrazy mają miarę zero
odpowiednio z Kroków~2, 3 i~4.
Przeliczalne pokrycia mapami przenoszą
wynik na $M$ i $N$. Zbiór miary zero
nie zawiera niepustego otwartego podzbioru
$N$, więc wartości regularne są gęste.
\end{proof}

\begin{uwaga}[Zakres twierdzenia]
Sard mówi o \emph{obrazach} punktów krytycznych,
nie o tym, że punktów krytycznych jest mało.
Funkcja stała ma wszystkie punkty krytyczne,
ale tylko jedną wartość krytyczną.
Dla odwzorowań klasy $C^r$ istnieje również
wersja przy założeniu całkowitoliczbowym
$r>\max\{m-n,0\}$; powyższy dowód dotyczy
wersji gładkiej, z której korzystamy dalej.
Jeśli dziedzina ma brzeg, stosujemy wersję bez brzegu
oddzielnie do $f|_{\operatorname{int}M}$ i
$f|_{\partial M}$. Suma ich zbiorów wartości krytycznych
ma miarę zero. Poza nią obie mapy są regularne na
odpowiednich przeciwobrazach; na brzegu surjektywność
$\dd f|_{T\partial M}$ pociąga także surjektywność
$\dd f$ na całym $TM$. Późniejsze twierdzenie o brzegu
wyjaśni, dlaczego potrzebny jest ten silniejszy warunek.
\end{uwaga}

\begin{wniosek}[Typowa poziomica]
\label{cor:sard-poziomica}
Dla prawie każdego $y\in N^n$ zbiór
$f^{-1}(y)$ jest pusty albo gładką
podrozmaitością wymiaru $m-n$, a w każdym
$x\in f^{-1}(y)$
$T_xf^{-1}(y)=\ker\dd f_x$.
\end{wniosek}
\begin{proof}
Poza zbiorem wartości krytycznych z
Twierdzenia~\ref{thm:sard} wartość $y$
jest regularna. Stosujemy
Twierdzenie~\ref{tw:przeciwobraz}; przypadek
pustego przeciwobrazu traktujemy osobno.
\end{proof}

\section{Transwersalność: precyzyjny sens przecięcia}
\label{sec:transwersalnosc}

Wartość regularna to szczególny przypadek
położenia ogólnego. Zamiast punktu $y\in N$
możemy badać trafienie odwzorowania w całą
podrozmaitość $Z\subset N$.

\begin{definicja}[Odwzorowanie transwersalne]
\index{odwzorowanie transwersalne@Odwzorowanie transwersalne}
\label{def:transwersalnosc}
Niech $f:M\to N$ będzie gładkie, a
$Z\subset N$ podrozmaitością zanurzoną w sensie
Rozdziału~\ref{ch:rozmaitosci}. Piszemy
$f\pitchfork Z$, gdy dla każdego
$x\in f^{-1}(Z)$ zachodzi
\begin{equation}\label{eq:transwersalnosc}
 \dd f_x(T_xM)+T_{f(x)}Z=T_{f(x)}N.
\end{equation}
Jeśli $f(M)\cap Z=\varnothing$, warunek jest spełniony
w sposób pusty. Dwie podrozmaitości $A,B\subset N$
przecinają się \emph{transwersalnie},
$A\pitchfork B$, jeśli
$T_xA+T_xB=T_xN$ przy każdym $x\in A\cap B$.
\end{definicja}

Suma w~\eqref{eq:transwersalnosc} jest sumą
\emph{podprzestrzeni}, a nie sumą wektorów
o ustalonych długościach. Jeśli $Z=\{y\}$,
to $T_yZ=0$ i otrzymujemy dokładnie warunek,
że $y$ jest wartością regularną $f$.
Gdy $f$ jest inkluzją $A\hookrightarrow N$,
otrzymujemy definicję transwersalnego
przecięcia $A$ z $Z$.

\begin{figure}[htbp]
\centering
\begin{tikzpicture}[font=\small,line cap=round]
 \begin{scope}[shift={(-2.6,0)}]
  \fill[manifoldblue!5] (-1.75,-1.3) rectangle (1.75,1.3);
  \draw[manifoldblue!85!black,very thick] (-1.45,0)--(1.45,0);
  \draw[manifoldgreen!75!black,very thick] (-1.06,-1.08)--(1.06,1.08);
  \fill[manifoldred] (0,0) circle (2.1pt);
  \node[manifoldblue!75!black] at (1.33,-.23) {$Z$};
  \node[manifoldgreen!65!black] at (.90,1.10) {$A$};
  \node at (0,-1.65) {przecięcie transwersalne};
 \end{scope}
 \begin{scope}[shift={(2.6,0)}]
  \fill[manifoldblue!5] (-1.75,-1.3) rectangle (1.75,1.3);
  \draw[manifoldblue!85!black,very thick] (-1.45,0)--(1.45,0);
  \draw[manifoldgold!80!black,very thick]
     plot[domain=-1.28:1.28,samples=55] (\x,{.63*\x*\x});
  \fill[manifoldred] (0,0) circle (2.1pt);
  \node[manifoldblue!75!black] at (1.33,-.23) {$Z$};
  \node[manifoldgold!75!black] at (1.25,1.12) {$A$};
  \node at (0,-1.65) {styczność: brak transwersalności};
 \end{scope}
\end{tikzpicture}
\caption{Po lewej dwie proste w płaszczyźnie mają
różne kierunki styczne, które razem rozpinają
$T_p\R^2$. Po prawej parabola dotyka osi:
obie przestrzenie styczne są tą samą prostą.}
\label{fig:transwersalne-przeciecia}
\end{figure}

\begin{twierdzenie}[Przeciwobraz podrozmaitości]
\label{thm:transwersalny-przeciwobraz}
Niech $Z\subset N^n$ będzie podrozmaitością
zanurzoną kowymiaru $c$, a $f:M^m\to N$
niech spełnia $f\pitchfork Z$. Wtedy
$f^{-1}(Z)$ jest pusty albo jest
podrozmaitością zanurzoną wymiaru $m-c$.
W punkcie $x\in f^{-1}(Z)$ zachodzi
\begin{equation}\label{eq:transwersalny-styczna}
 T_xf^{-1}(Z)
  =(\dd f_x)^{-1}\bigl(T_{f(x)}Z\bigr).
\end{equation}
Jeśli $m<c$, przeciwobraz jest pusty.
\end{twierdzenie}
\begin{proof}
Wokół $z=f(x)$ wybierzmy mapę
$\psi:V\to\R^{n-c}\times\R^c$, w której
$\psi(Z\cap V)=\psi(V)\cap(\R^{n-c}\times\{0\})$.
Niech $h=\operatorname{pr}_{\R^c}\circ\psi$.
Wtedy lokalnie $Z=h^{-1}(0)$ oraz
$\ker\dd h_z=T_zZ$.
Reguła łańcucha pokazuje, że
$\dd(h\circ f)_x=\dd h_z\circ\dd f_x$.
Ta różniczka jest surjekcją na $\R^c$
dokładnie wtedy, gdy
$\dd f_x(T_xM)+\ker\dd h_z=T_zN$;
jest to~\eqref{eq:transwersalnosc}.
Zatem $0$ jest lokalnie wartością regularną
$h\circ f$. Twierdzenie~\ref{tw:przeciwobraz}
daje podrozmaitość kowymiaru $c$ i
\[
 \ker\dd(h\circ f)_x
  =\{v:\dd f_xv\in\ker\dd h_z\},
\]
co jest~\eqref{eq:transwersalny-styczna}.
Jeśli $m<c$, surjekcja $\R^m\to\R^c$
nie istnieje. Nie może więc istnieć punkt
przeciwobrazu przy założonej transwersalności.
\end{proof}

\begin{przyklad}[Prosta i parabola]
W $N=\R^2$ niech $Z=\{(x,0)\}$.
Dla $f(t)=(t,t)$ mamy
$f^{-1}(Z)=\{0\}$ oraz
$\dd f_0(\R)=\operatorname{span}(1,1)$,
który wraz z $T_{(0,0)}Z=\operatorname{span}(1,0)$
rozpina $\R^2$. Dla $g(t)=(t,t^2)$
przeciwobraz także jest pojedynczym punktem,
ale $\dd g_0(\R)=T_{(0,0)}Z$:
sam fakt, że punktów przecięcia jest skończenie
wiele, \emph{nie} oznacza transwersalności.
Na Rysunku~\ref{fig:transwersalne-przeciecia}
przedstawiono oba przypadki.
\end{przyklad}

\begin{twierdzenie}[Przecięcie i iloczyn włóknisty]
\label{thm:transwersalne-przeciecie}
Jeżeli podrozmaitości $A^a,B^b\subset N^n$
spełniają $A\pitchfork B$, to $A\cap B$
jest pusty albo jest podrozmaitością wymiaru $a+b-n$ oraz
\[
 T_x(A\cap B)=T_xA\cap T_xB.
\]
Ogólniej, dla gładkich $f:M^m\to N^n$ i
$g:P^p\to N^n$ załóżmy, że przy każdym
$f(x)=g(y)=z$ mamy
\begin{equation}\label{eq:transwersalne-dwa-odwz}
 \dd f_x(T_xM)+\dd g_y(T_yP)=T_zN.
\end{equation}
Wówczas
$M\times_N P=\{(x,y):f(x)=g(y)\}$
jest pusty albo jest podrozmaitością $M\times P$ wymiaru
$m+p-n$, a jej przestrzeń styczna składa się z
par $(u,v)$ spełniających $\dd f_xu=\dd g_yv$.
\end{twierdzenie}
\begin{proof}
Niech $\Delta_N=\{(z,z):z\in N\}$ będzie diagonalą
w $N\times N$. Jej przestrzeń styczna w $(z,z)$
to $\{(w,w):w\in T_zN\}$. Obliczamy
\[
 T_{(z,z)}(N\times N)/T_{(z,z)}\Delta_N
       \cong T_zN,\qquad [(u,v)]\longmapsto u-v.
\]
Odwzorowanie $(f,g):M\times P\to N\times N$
jest transwersalne do diagonali dokładnie wtedy,
gdy $(u,v)\mapsto\dd f_xu-\dd g_yv$ jest surjekcją,
czyli gdy zachodzi~\eqref{eq:transwersalne-dwa-odwz}.
Twierdzenie~\ref{thm:transwersalny-przeciwobraz}
daje podrozmaitość kowymiaru $n$.
Jej przestrzeń styczna jest jądrem powyższej
różnicy, skąd oba wzory. Dla inkluzji
$A\hookrightarrow N$ i $B\hookrightarrow N$
iloczyn włóknisty utożsamiamy z $A\cap B$
przez $(z,z)\mapsto z$. Warunek staje się
$T_zA+T_zB=T_zN$, a jądro jest przecięciem
obu przestrzeni stycznych.
\end{proof}

\begin{uwaga}[Oczekiwany wymiar]
Liczba $a+b-n$ to \emph{oczekiwany wymiar}
przecięcia. Ujemny wynik wymusza pustkę przy
transwersalności. Wynik zero oznacza izolowane
punkty przecięcia. Jeśli $A$ jest zwarta, a $B$ domknięta
(w szczególności gdy obie są zwarte), przecięcie jest
zwarte i dyskretne, więc skończone. Sama dyskretność
bez zwartości nie wyklucza nieskończenie wielu punktów.
To właśnie przypadek, w którym później można
liczyć przecięcia ze znakami, używając orientacji
z Sekcji~\ref{sec:formy-orientacja}.
\end{uwaga}

\section{Brzeg i kontrola miejsca przecięcia}
\label{sec:transwersalnosc-brzeg}

W konstrukcjach topologicznych pojawiają się
przestrzenie z brzegiem. Sam warunek
$f\pitchfork Z$ we wnętrzu nie gwarantuje
poprawnego zachowania przeciwobrazu na brzegu.

\begin{twierdzenie}[Przeciwobraz transwersalny z brzegiem]
\label{thm:transwersalnosc-brzeg}
Niech $M^m$ będzie rozmaitością z brzegiem,
$N$ bez brzegu, $Z\subset N$ podrozmaitością
kowymiaru $c$. Jeżeli zarówno $f:M\to N$,
jak i $f|_{\partial M}:\partial M\to N$
są transwersalne do $Z$, to
$W=f^{-1}(Z)$ jest pusty albo jest rozmaitością z brzegiem
wymiaru $m-c$ i
\[
 \partial W=W\cap\partial M.
\]
\end{twierdzenie}
\begin{proof}
We wnętrzu $M$ stosujemy
Twierdzenie~\ref{thm:transwersalny-przeciwobraz}.
Przy $x\in\partial M\cap W$ wybierzmy współrzędne
półprzestrzeni $(r,u_2,\ldots,u_m)$,
$r\geq0$, oraz lokalną funkcję definiującą
$Z$, $h:N\to\R^c$ jak w poprzednim dowodzie.
Połóżmy $H=h\circ f$. Transwersalność
$f|_{\partial M}$ oznacza, że
$\dd H_x:T_x\partial M\to\R^c$ jest surjekcją.
Możemy więc wybrać $c$ kierunków
\emph{stycznych do brzegu}, na których
współrzędne $H$ tworzą niezależny układ.
Po przestawieniu zmiennych $u$ odpowiedni minor $c\times c$
jest odwracalny. Gładkość na półprzestrzeni pozwala
przedłużyć $H$ lokalnie przez $r=0$. Twierdzenie o funkcji
odwrotnej dla mapy $(r,u)\mapsto(r,H(r,u),v)$ stosujemy
w tym otwartym otoczeniu, a potem ograniczamy do $r\geq0$.
Jej różniczka jest odwracalna właśnie dzięki wybranemu
minorowi i zachowaniu pozostałych współrzędnych.
Ponieważ $r$ pozostało współrzędną brzegową,
zbiór $H=0$ ma w tych współrzędnych postać
$\{(r,0,v):r\geq0\}$. Jego brzeg jest dokładnie
$\{r=0,H=0\}=W\cap\partial M$.
Pozostałe $m-1-c$ współrzędne $v$ wraz z $r$
dają wymiar $m-c$. Przy $m<c$ cały przeciwobraz jest
pusty; przy $m=c$ nie może spotykać brzegu, ponieważ
surjekcja z jego $(m-1)$-wymiarowej przestrzeni stycznej
na $\R^c$ nie istnieje.
\end{proof}

\begin{przyklad}[Dlaczego warunek na brzegu jest potrzebny]
Weźmy górną półpłaszczyznę
$M=\{(x,y)\in\R^2:y\geq0\}$,
$N=\R$ i $Z=\{0\}$. Funkcja $f(x,y)=y$
jest submersją w każdym punkcie $M$,
lecz jej ograniczenie do brzegu jest stale
równe zero. Mamy $f^{-1}(0)=\partial M$,
czyli prostą \emph{bez} brzegu; równość
$\partial f^{-1}(Z)=f^{-1}(Z)\cap\partial M$
byłaby tutaj fałszywa. W praktyce należy
sprawdzać oba warunki osobno.
\end{przyklad}

\section{Transwersalność z parametrem}
\label{sec:transwersalnosc-param}

Nie należy rozumieć „ogólnego położenia” jako
założenia o losowości rysunku. Chodzi o precyzyjne
twierdzenie: w gładkiej rodzinie dopuszczalnych
przesunięć złe parametry tworzą zbiór miary zero.

\begin{twierdzenie}[Transwersalność parametryczna]
\label{thm:transwersalnosc-param}
Niech $M,A,N$ będą gładkimi rozmaitościami bez brzegu,
$Z\subset N$ podrozmaitością zanurzoną, a
$F:M\times A\to N$ gładkim odwzorowaniem.
Jeżeli $F\pitchfork Z$, to dla prawie każdego
$a\in A$ odwzorowanie $F_a(x)=F(x,a)$ spełnia
$F_a\pitchfork Z$. W szczególności takie
parametry są gęste w $A$.
\end{twierdzenie}
\begin{proof}
Ponieważ $F\pitchfork Z$,
$W=F^{-1}(Z)$ jest podrozmaitością
$M\times A$ na mocy
Twierdzenia~\ref{thm:transwersalny-przeciwobraz}.
Rozważmy rzutowanie $\pi:W\to A$.
W punkcie $(x,a)\in W$, $z=F(x,a)$,
\[
 T_{(x,a)}W
 =\{(u,b):\dd_xF_{(x,a)}u
              +\dd_aF_{(x,a)}b\in T_zZ\}.
\]
Zatem $\dd\pi_{(x,a)}$ jest surjekcją
dokładnie wtedy, gdy dla każdego $b\in T_aA$
można dobrać $u\in T_xM$ tak, aby
$\dd_aF(b)+\dd_xF(u)\in T_zZ$.
Inaczej mówiąc,
\[
 \dd_aF(T_aA)\subset
                 \dd_xF(T_xM)+T_zZ.
\]
Łącznie z założeniem
$\dd_xF(T_xM)+\dd_aF(T_aA)+T_zZ=T_zN$
jest to równoważne
$\dd_xF(T_xM)+T_zZ=T_zN$, czyli
transwersalności $F_a$ w punkcie $x$.
Wartość $a$ jest regularna dla $\pi$
wtedy i tylko wtedy, gdy $F_a$ jest
transwersalne do $Z$ we wszystkich punktach
przeciwobrazu. Sard zastosowany do
$\pi:W\to A$ kończy dowód, jeśli $\dim A\geq1$.
Dla $\dim A=0$ przestrzeń parametrów jest dyskretna,
a $T_aA=0$. Wtedy $\dd F=\dd_xF$ na każdej składowej
$M\times\{a\}$ i założenie daje tezę dla każdego $a$,
bez odwoływania się do Sarda.
\end{proof}

\begin{przyklad}[Rozszczepienie stycznego przecięcia]
\label{prz:transwersalne-parabola}
Niech $M=N=\R$, $Z=\{0\}$ oraz
$F(t,a)=t^2-a$. Pochodna w kierunku parametru
wynosi $\partial_aF=-1$, więc cała rodzina
$F:\R^2\to\R$ jest submersją.
Dla $a>0$ mapa $F_a$ ma dwa regularne zera
$t=\pm\sqrt a$, bo $F_a'(t)=2t\ne0$.
Dla $a<0$ nie ma zer, a $a=0$ daje
styczne, nieregularne zero $t=0$.
Jedyny zły parametr to $0$.
\end{przyklad}

\begin{figure}[htbp]
\centering
\begin{tikzpicture}[font=\small,line cap=round]
 \foreach \panel/\aval/\lab in {-3.15/-.5/$a=-\frac12$,0/0/$a=0$,3.15/.5/$a=\frac12$}{
  \begin{scope}[shift={(\panel,0)}]
   \draw[->,manifoldblue!65!black] (-1.18,0)--(1.22,0);
   \draw[->,manifoldblue!65!black] (0,-.74)--(0,1.80);
   \draw[manifoldgold!75!black,very thick]
    plot[domain=-1.05:1.05,samples=60,variable=\t]
       (\t,{(\t)^2-(\aval)});
   \node at (0,-1.02) {\lab};
  \end{scope}
 }
 \fill[manifoldred] (0,0) circle (2pt);
 \fill[manifoldred] ({3.15-sqrt(.5)},0) circle (2pt);
 \fill[manifoldred] ({3.15+sqrt(.5)},0) circle (2pt);
\end{tikzpicture}
\caption{Poziom zero rodziny $F_a(t)=t^2-a$.
Zmiana parametru rozbija styczność w dwa
przecięcia transwersalne albo usuwa przecięcie;
sam parametr $a=0$ jest wyjątkowy.}
\label{fig:transwersalnosc-perturbacja}
\end{figure}

Twierdzenie parametryczne nie mówi jeszcze, że
\emph{każde} konkretne odwzorowanie ma dostatecznie
bogatą rodzinę perturbacji. Przy zwartej dziedzinie
możemy taką rodzinę zbudować bez odwoływania się
do twierdzenia Whitneya o zanurzeniach.

Wyjaśnijmy używaną tu topologię przybliżania. Wybierzmy
skończenie wiele map $x_i$ w $M$ i $y_i$ w $N$ oraz
zwartych zbiorów $K_i$ w dziedzinach $x_i$ tak, by ich
wnętrza pokrywały $M$, a $f(K_i)$ leżały w dziedzinach $y_i$.
Zapewniają to ciągłość $f$ i zwartość $M$.
Dla odwzorowań $g$ bliskich $f$ wymagamy
$g(K_i)$ w tych samych mapach i kontrolujemy liczby
\[
 \max_i\max_{|\alpha|\leq r}\sup_{u\in x_i(K_i)}
 \left|\partial^\alpha(y_i\circ g\circ x_i^{-1})(u)
       -\partial^\alpha(y_i\circ f\circ x_i^{-1})(u)\right|.
\]
Małość tych liczb określa topologię $C^r$. Reguła łańcucha
i zwartość pokazują, że zmiana skończonego układu map
daje tę samą topologię. Topologia $C^\infty$ jest generowana
przez te warunki dla wszystkich skończonych $r$:
zbieżność oznacza zbieżność w każdym $C^r$, a podstawowe
otoczenie zadaje tylko skończenie wiele takich ograniczeń.
Nie wymagamy jednego oszacowania jednocześnie dla
pochodnych wszystkich rzędów.

\begin{twierdzenie}[Przybliżanie odwzorowań transwersalnymi]
\label{thm:gestosc-transwersalnych}
Niech $M$ będzie zwartą rozmaitością bez brzegu,
$N$ rozmaitością, a $Z\subset N$ podrozmaitością
zanurzoną. Każde gładkie $f:M\to N$ można
dowolnie blisko przybliżyć w topologii $C^\infty$
gładkim $g:M\to N$ transwersalnym do $Z$.
Można przy tym wybrać gładką homotopię
od $f$ do $g$.
Jeśli dodatkowo $Z$ jest domknięta w $N$,
transwersalność jest otwarta w topologii $C^1$
na gładkich odwzorowaniach z $M$ do $N$.
\end{twierdzenie}
\begin{proof}
\emph{Budowa skończonej liczby kierunków perturbacji.}
Wybierzmy na $N$ dowolną metrykę riemannowską.
Przeciągnięta wiązka $f^*TN$ ma nad małymi
otoczeniami $U_i\subset M$ gładkie ramy.
Ponieważ $M$ jest zwarta, możemy wybrać skończenie
wiele mniejszych $V_i\Subset U_i$, które pokrywają
$M$. Dla każdego $i$ wybierzmy funkcję odcinającą
$\rho_i$, równą $1$ na $V_i$ i o nośniku w $U_i$.
Pomnóżmy każdy wektor lokalnej ramy $f^*TN|_{U_i}$
przez $\rho_i$ i przedłużmy przez zero. Otrzymujemy
skończenie wiele globalnych przekrojów
$s_1,\ldots,s_k$ wiązki $f^*TN$ takich, że
ich wartości rozpinają $T_{f(x)}N$ w każdym $x$.

Niech $a=(a_1,\ldots,a_k)$ przebiega małą kulę
w $\R^k$. Definiujemy
\[
 F(x,a)=\exp^N_{f(x)}
              \left(\sum_{j=1}^k a_js_j(x)\right).
\]
Metryka nie musi być zupełna: wykładnicze
odwzorowanie jest określone przy wektorze zero,
a zwartość $M$ pozwala wziąć jedną wystarczająco
małą kulę parametrów dla wszystkich $x$.
W punkcie $a=0$ pochodna $F$ względem parametrów
przyjmuje na standardowym kierunku $j$ wartość
$s_j(x)$, ponieważ $\dd(\exp^N_{f(x)})_0=\id$.
Jest więc surjektywna na $T_{f(x)}N$.
Surjektywność jest warunkiem otwartym;
przez zwartość $M$ pozostaje prawdziwa na
$M\times A$ po zmniejszeniu kuli $A$.
Cała mapa $F$ jest zatem submersją,
w szczególności $F\pitchfork Z$.
Twierdzenie~\ref{thm:transwersalnosc-param}
dostarcza parametrów $a$ transwersalnych do
$Z$ dowolnie blisko zera. Dla
$g=F_a$ mamy $F_a\to F_0=f$ w $C^\infty$:
w ustalonych mapach każda pochodna po $x$ zależy
ciągle od $(x,a)$, a na skończonym pokryciu zwartych
$K_i$ zbieżność przy $a\to0$ jest jednostajna dla każdego
skończonego rzędu. Ścieżka $t\mapsto F_{ta}$ daje
gładką homotopię, ponieważ kula $A$ jest wypukła.

\emph{Otwartość przy domkniętym $Z$.}
Załóżmy, że $f\pitchfork Z$, lecz pewne
$f_j\to f$ w $C^1$ nie są transwersalne.
Wybierzmy punkty $x_j$ naruszenia warunku.
Po przejściu do podciągu $x_j\to x$ dzięki
zwartości $M$. Ponieważ $Z$ jest domknięta,
$f(x)=\lim f_j(x_j)$ należy do $Z$.
W lokalnej mapie spłaszczającej $Z$
transwersalność oznacza, że ustalony
minor macierzy pochodnej normalnych
współrzędnych jest niezerowy. Dla $f$
jest taki minor w $x$; ciągłość i zbieżność
$C^1$ dają ten sam niezerowy minor dla
$f_j$ w $x_j$ przy dużym $j$.
Otrzymujemy sprzeczność.
\end{proof}

\begin{wniosek}[Perturbacja względna]
\label{cor:transwersalnosc-wzgledna}
Załóżmy dodatkowo w poprzednim twierdzeniu,
że $Z$ jest domknięta i $f\pitchfork Z$
na otoczeniu domkniętego $C\subset M$.
Wówczas $g$ można wybrać równe $f$
na pewnym otoczeniu $C$.
\end{wniosek}
\begin{proof}
Wybierzmy otoczenia
$C\subset U_0\Subset U_1\Subset U$,
gdzie $f\pitchfork Z$ na $U$.
Powtórzmy konstrukcję przekrojów z poprzedniego
dowodu tylko nad zwartym zbiorem
$M\setminus U_1$, dobierając ich nośniki
w $M\setminus\overline U_0$.
Otrzymana rodzina $F_a$ jest identyczna z
$f$ na $U_0$, a jej pochodna parametryczna
rozpina $TN$ na $M\setminus U_1$.
Na zwartym $\overline U_1\subset U$ dostatecznie małe
perturbacje zachowują transwersalność $f$. Stosujemy
tu dowód otwartości na zwartym podzbiorze, nie twierdzenie
o rozmaitości $\overline U_1$: hipotetyczne punkty
naruszenia miałyby podciąg z granicą w $\overline U_1$,
a ten sam argument z niezerowym minorem dawałby sprzeczność.
Stąd cała rodzina $F:M\times A\to N$
jest transwersalna do $Z$.
Wybieramy parametr z
Twierdzenia~\ref{thm:transwersalnosc-param}
dowolnie blisko zera. Na $U_0$ wynikowe
$g=F_a$ jest dokładnie $f$. Cała homotopia $F_{ta}$
również pozostaje tam równa $f$.
\end{proof}

Jeśli $M$ ma brzeg, a rodzina $F$ i jej
ograniczenie $F|_{\partial M\times A}$
są transwersalne do $Z$, można zastosować
Sarda do dwóch map bez brzegu:
$\pi|_{\operatorname{int}W}$ oraz $\pi|_{\partial W}$,
gdzie $W=F^{-1}(Z)$ i $\partial W=W\cap(\partial M\times A)$.
Pierwsza kontroluje transwersalność $F_a$ we wnętrzu,
druga transwersalność jego ograniczenia do brzegu.
Ta ostatnia pociąga też transwersalność $F_a$ na całej
przestrzeni stycznej $T_xM$ w punkcie brzegowym.
Suma dwóch zbiorów parametrów wyjątkowych
nadal ma miarę zero. Dla pozostałych
parametrów otrzymujemy oba warunki
Twierdzenia~\ref{thm:transwersalnosc-brzeg}.

\section{Zera przekrojów i punkty krytyczne funkcji}
\label{sec:transwersalnosc-przekroje}

Nie każdy problem ma postać przecięcia dwóch
podzbiorów jednej przestrzeni. Często szukamy
zer pola wektorowego, $1$-formy albo innego
przekroju wiązki. Wtedy transwersalność działa
względem \emph{przekroju zerowego}.

\begin{twierdzenie}[Regularne zera przekroju]
\label{thm:zera-przekroju}
Niech $E\to M^m$ będzie wiązką wektorową
rzędu $r$, a $s:M\to E$ gładkim przekrojem.
W lokalnej trywializacji $s(x)=(x,\sigma(x))$.
W punkcie $s(x)=0$ pochodna
$\D_xs:=\dd\sigma_x:T_xM\to E_x$
nie zależy od wyboru trywializacji
po naturalnym utożsamieniu włókien.
Jeżeli $\D_xs$ jest surjektywna przy
każdym zerze, to zbiór zer $s^{-1}(0_E)$
jest pusty albo jest podrozmaitością kowymiaru $r$, a
$T_xs^{-1}(0_E)=\ker\D_xs$.
\end{twierdzenie}
\begin{proof}
W trywializacji przekrój zerowy ma postać
$\{(x,0)\}$, a kierunki normalne do niego
odpowiadają współrzędnym włókna.
Różniczka $s(x)=(x,\sigma(x))$ ma postać
$v\mapsto(v,\dd\sigma_xv)$. Warunek
$s\pitchfork 0_E$ jest więc dokładnie
surjektywnością $\dd\sigma_x$.
Twierdzenie~\ref{thm:transwersalny-przeciwobraz}
daje wymiar i przestrzeń styczną.
Przy zmianie trywializacji
$\widetilde\sigma(y)=A(y)\sigma(y)$,
gdzie $A(y)$ jest odwracalną macierzą.
W punkcie z $\sigma(x)=0$ reguła iloczynu
daje $\dd\widetilde\sigma_x=A(x)\dd\sigma_x$:
wyraz $(\dd A_x)\sigma(x)$ znika.
Surjektywność i jądro nie zależą więc
od trywializacji. W sposób wewnętrzny w zerze $s(x)$
mamy kanoniczny rozkład
$T_{0_x}E\cong T_xM\oplus E_x$: pierwsza część jest
styczna do przekroju zerowego, a druga jest pionowa.
Pionowa składowa $\dd s_x$ to właśnie $\D_xs$.
Dla $r>m$ surjekcja jest niemożliwa, więc przekrój
transwersalny nie ma zer.
\end{proof}

\begin{przyklad}[Zera pola na płaszczyźnie]
Pole $X(x,y)=x\,\partial_x+y\,\partial_y$ ma
jedno zero, a jego pochodna pionowa w zerze
jest macierzą jednostkową: zero jest
transwersalne. Pole
$Y(x,y)=x^2\partial_x+y\,\partial_y$
też ma jedno zero, lecz macierz
$\D_0Y=\operatorname{diag}(0,1)$
nie jest surjektywna. Sama izolacja zera
nie wystarcza do transwersalności.
\end{przyklad}

Dla funkcji $f:M\to\R$ przekrojem, który
nas interesuje, jest $df:M\to T^*M$.
To pozwala wyjaśnić, skąd w przyszłej
teorii Morse'a pojawia się hesjan.

\begin{definicja}[Hesjan i niezdegenerowany punkt krytyczny]
\index{hesjan i niezdegenerowany punkt krytyczny@Hesjan i niezdegenerowany punkt krytyczny}
\label{def:morse-pochodna}
Jeśli $\dd f_p=0$, to w dowolnej mapie
$(x^1,\ldots,x^m)$ wokół $p$ definiujemy
\[
 \operatorname{Hess}_p f(v,w)
 =\sum_{i,j}
    \frac{\partial^2 f}{\partial x^i\partial x^j}(p)
                   v^iw^j.
\]
Punkt $p$ nazywamy \emph{niezdegenerowanym},
gdy ta forma dwuliniowa ma nieosobliwą macierz.
Funkcja, której wszystkie punkty krytyczne
są niezdegenerowane, nazywa się
\emph{funkcją Morse'a}.
\end{definicja}

Hesjan w punkcie krytycznym jest niezależny
od mapy. Druga pochodna złożenia
$f\circ\varphi^{-1}$ z przejściem współrzędnych
ma wyraz z drugą pochodną $f$ przeniesiony
przez jakobian oraz wyraz z pierwszą pochodną
$f$ pomnożoną przez drugie pochodne przejścia.
Ten ostatni wyraz znika właśnie dlatego,
że $\dd f_p=0$. Dokładniej, jeśli $x=x(y)$ jest zmianą map,
to w punkcie krytycznym
\[
 \frac{\partial^2 f}{\partial y^a\partial y^b}(p)
 =\sum_{i,j}\frac{\partial x^i}{\partial y^a}(p)
             \frac{\partial x^j}{\partial y^b}(p)
             \frac{\partial^2 f}{\partial x^i\partial x^j}(p).
\]
Jest to prawo transformacji formy dwuliniowej,
$H_y=J^{\mathsf T}H_xJ$, więc niezdegenerowanie jest
niezależne od mapy. Po wyborze dowolnej koneksji
Leviego-Civity wcześniejszy hesjan ma współrzędne
$(\nabla df)_{ij}=\partial_i\partial_jf-\Gamma^k_{ij}\partial_kf$;
w punkcie krytycznym daje tę samą formę. Nie wymaga
ona zatem wyboru metryki w tym punkcie.
W mapie wiązki kostycznej
przekrój $df$ ma współrzędne
$(x,(\partial_1f,\ldots,\partial_m f))$.
Jego pochodna pionowa w punkcie krytycznym
jest więc macierzą
$(\partial_i\partial_j f)(p)$. Dostajemy
równoważność
\begin{equation}\label{eq:morse-transwersalnosc}
 df\pitchfork 0_{T^*M}
 \quad\Longleftrightarrow\quad
 \text{każdy punkt krytyczny $f$ jest
 niezdegenerowany.}
\end{equation}

\begin{twierdzenie}[Istnienie i typowość funkcji Morse'a]
\label{thm:gestosc-morse}
Na każdej zwartej rozmaitości gładkiej
bez brzegu funkcje Morse'a są gęste wśród
funkcji gładkich w topologii $C^\infty$
i tworzą zbiór otwarty w topologii $C^2$.
Każda taka funkcja ma skończenie wiele
punktów krytycznych. Można dodatkowo
osiągnąć, by ich wartości były parami różne.
\end{twierdzenie}
\begin{proof}
\emph{Gęstość.}
Weźmy skończone mapy $(U_i,x_i^1,\ldots,x_i^m)$
i funkcje odcinające $\rho_i$ tak, aby mniejsze
obszary, na których $\rho_i=1$, pokrywały $M$,
a $\operatorname{supp}\rho_i\Subset U_i$.
Funkcje $\rho_i x_i^j$, przedłużone przez
zero poza $U_i$, są globalnie gładkie.
Niech
\[
 f_a=f+\sum_{i,j}a_{ij}\rho_i x_i^j.
\]
W dowolnym $x$ leżącym w mniejszym obszarze
mapy $i$ mamy
$\dd(\rho_i x_i^j)_x=\dd x_i^j$.
Kowektory te tworzą bazę $T_x^*M$.
Zatem parametryczna rodzina przekrojów
$(x,a)\mapsto df_a(x)$ ma surjektywną
pochodną pionową w kierunkach $a$:
jest transwersalna do przekroju zerowego
$0_{T^*M}$. Twierdzenie
~\ref{thm:transwersalnosc-param}
zastosowane do tej rodziny mówi, że
dla prawie każdego małego $a$ przekrój
$df_a$ jest transwersalny do zera.
Równoważność~\eqref{eq:morse-transwersalnosc}
czyni $f_a$ funkcją Morse'a. Ponieważ
$a\to0$ można wybrać dowolnie,
$f_a\to f$ w $C^\infty$.

\emph{Otwartość i skończoność.}
Jeśli $f$ jest Morse'a, przekrój $df$
jest transwersalny do domkniętego
przekroju zerowego. Dowód otwartości
z Twierdzenia~\ref{thm:gestosc-transwersalnych}
stosuje się także do przekrojów;
gdy $g$ jest bliskie $f$ w $C^2$,
$dg$ jest bliskie $df$ w $C^1$.
To daje otwartość. Zbiór krytyczny
jest zerowymiarową podrozmaitością
na mocy Twierdzenia~\ref{thm:zera-przekroju},
więc jego punkty są izolowane.
Jest także zamknięty w zwartym $M$;
zatem jest skończony.

\emph{Różne wartości.}
Niech krytyczne punkty funkcji Morse'a
będą $p_1,\ldots,p_\ell$. Wybierzmy
rozłączne małe otoczenia i funkcje
odcinające $\beta_i$ równe $1$ w jeszcze
mniejszych otoczeniach $p_i$.
Zastąpmy $f$ przez
$f+\sum_i c_i\beta_i$.
Przy dostatecznie małych $c_i$ w obszarach,
gdzie zmieniają się $\beta_i$, nie pojawiają
się nowe punkty krytyczne: na zwartym
zbiorze tych obszarów $|\dd f|$ ma dodatnie
minimum, a $|\sum_i c_i\dd\beta_i|$ może
być od niego mniejsze. W pobliżu $p_i$
nowa funkcja różni się od $f$ o stałą
$c_i$, więc zachowuje punkty i hesjany.
Nowe wartości to $f(p_i)+c_i$; wybieramy
małe $c_i$ poza skończoną sumą hiperpłaszczyzn
$f(p_i)+c_i=f(p_j)+c_j$.
\end{proof}

\begin{przyklad}[Mała perturbacja funkcji]
Na $\R^2$ funkcja $f(x,y)=x^2$ ma całą prostą
$x=0$ punktów krytycznych; hesjan
$\operatorname{diag}(2,0)$ jest osobliwy.
Po dodaniu $\varepsilon y^2$ (perturbacji małej wraz
z pochodnymi na każdym ustalonym zbiorze zwartym) dla
$\varepsilon\ne0$ pozostaje tylko punkt
$(0,0)$, a hesjan
$\operatorname{diag}(2,2\varepsilon)$
jest odwracalny. Gdy $\varepsilon>0$
otrzymujemy minimum, a gdy
$\varepsilon<0$ punkt siodłowy.
To lokalny obraz perturbacji; twierdzenie
poprzednie zapewnia analogiczne
rozwiązanie globalne na zwartej rozmaitości.
\end{przyklad}

\section{Dlaczego transwersalność pojawi się u Smale'a?}
\label{sec:transwersalnosc-smale}

Niech $(M^n,g)$ będzie rozmaitością riemannowską,
a $f:M\to\R$ funkcją Morse'a,
a $X=-\nabla_g f$ jej polem spadku. W punkcie
krytycznym $p$ \emph{indeks Morse'a}
$\lambda(p)$ to liczba ujemnych wartości
własnych formy $\operatorname{Hess}_p f$
po podniesieniu jednego indeksu metryką $g$,
liczonych z krotnościami. Prawo bezwładności Sylvestera
pokazuje, że ta liczba zależy od formy hesjanu,
a nie od pomocniczej metryki użytej do wyznaczenia
wartości własnych. W dalszej teorii
przepływu gradientowego skonstruujemy
rozmaitość niestabilną $W^u(p)$, złożoną
z punktów dążących do $p$ przy czasie
$t\to-\infty$, oraz stabilną $W^s(q)$,
złożoną z punktów dążących do $q$
przy $t\to+\infty$. W tych definicjach wymagamy
istnienia trajektorii przez całą odpowiednią półprostą
czasu. Na zwartej rozmaitości jest to automatyczne;
gdy pole nie jest zupełne, nie wolno tego pomijać.
Twierdzenie o
rozmaitościach stabilnych daje
\[
 \dim W^u(p)=\lambda(p),\qquad
 \dim W^s(q)=n-\lambda(q).
\]
Nie dowodzimy tutaj tego twierdzenia o równaniach
różniczkowych; poniżej wyprowadzamy już teraz
\emph{konsekwencję transwersalności} tych obiektów.
Globalnie traktujemy je jako rozmaitości z immersjami
do $M$; lokalnie każda wybrana gałąź jest podrozmaitością
z mapą plasterkową. Wniosek poniżej dotyczy takiego
lokalnego przecięcia, co nie wymaga z góry twierdzenia
o globalnym zanurzeniu obu zbiorów.

\begin{przyklad}[Siodło i dwa kierunki przepływu]
\label{prz:siodlo-transwersalne}
W $\R^2$ dla metryki euklidesowej i
$f(x,y)=x^2-y^2$ pole spadku wynosi
$X=(-2x,2y)$. Rozwiązanie startujące z
$(x_0,y_0)$ jest jawne:
\[
 x(t)=x_0e^{-2t},\qquad y(t)=y_0e^{2t}.
\]
Do $(0,0)$ przy $t\to+\infty$ dążą dokładnie
punkty z $y_0=0$, więc
$W^s(0)=\{y=0\}$. Przy $t\to-\infty$
dążą punkty z $x_0=0$, więc
$W^u(0)=\{x=0\}$. Ich styczne w siodle
rozpinają $\R^2$, jak przewiduje
nieosobliwy hesjan $\operatorname{diag}(2,-2)$.
\end{przyklad}

\begin{figure}[htbp]
\centering
\begin{tikzpicture}[>=Latex,font=\small,line cap=round,scale=1.35]
 \fill[manifoldblue!5] (-2,-1.68) rectangle (2,1.68);
 \draw[->,manifoldblue!80!black,very thick] (-1.72,0)--(-.19,0);
 \draw[->,manifoldblue!80!black,very thick] (1.72,0)--(.19,0);
 \draw[->,manifoldgold!80!black,very thick] (0,.19)--(0,1.50);
 \draw[->,manifoldgold!80!black,very thick] (0,-.19)--(0,-1.50);
 \draw[manifoldgreen!70!black,thin,
       postaction={decorate,decoration={markings,
              mark=at position .58 with {\arrow{Latex}}}}]
  plot[domain=1.46:.21,samples=40] (\x,{.24/\x});
 \fill[manifoldred] (0,0) circle (2.6pt)
    node[below right=3pt] {$p$};
 \node[manifoldblue!75!black,anchor=west] at (1.08,-.32)
       {$W^s(p)$};
 \node[manifoldgold!65!black,anchor=west] at (.15,1.27)
       {$W^u(p)$};
\end{tikzpicture}
\caption{Lokalny model siodła $f=x^2-y^2$.
Niebieska oś jest stabilna: przepływ kieruje się
ku $p$. Złota oś jest niestabilna: przepływ
oddala się od $p$. Zielona trajektoria ilustruje
ruch poza osiami; osie przecinają się transwersalnie.}
\label{fig:smale-siodlo}
\end{figure}

\begin{wniosek}[Wymiar przecięcia trajektorii]
\label{cor:smale-wymiar}
Jeżeli $W^u(p)$ i $W^s(q)$ przecinają się
transwersalnie, to ich przecięcie jest lokalnie
rozmaitością wymiaru
$\lambda(p)-\lambda(q)$.
Jeśli $p\ne q$, w każdym punkcie należącym
do łączącej je trajektorii niezerowe pole
$X$ jest styczne do obu rozmaitości.
Po lokalnym podzieleniu przez przesunięcie
czasu otrzymujemy wymiar
$\lambda(p)-\lambda(q)-1$.
\end{wniosek}
\begin{proof}
Podstawiamy podane wymiary do
Twierdzenia~\ref{thm:transwersalne-przeciecie}, stosowanego
do lokalnych gałęzi (lub do iloczynu włóknistego immersji):
$\lambda(p)+(n-\lambda(q))-n
=\lambda(p)-\lambda(q)$.
Obie rozmaitości są niezmiennicze
pod przepływem $X$, więc wektor prędkości
trajektorii należy do obu przestrzeni
stycznych. Dla połączenia różnych
punktów krytycznych nie znika on:
gdyby trajektoria dotarła do punktu
równowagi w skończonym czasie,
jednoznaczność równania przepływu
czyniłaby ją stałą. Twierdzenie o prostowaniu
pola~\ref{thm:flowbox}, zastosowane do ograniczenia $X$
na przecięciu, daje lokalną współrzędną czasu; usunięcie tej jednej współrzędnej
zmniejsza wymiar o jeden.
\end{proof}

W szczególności izolowanych trajektorii
po lokalnym podzieleniu przez czas oczekujemy przy
$\lambda(p)-\lambda(q)=1$. Gdy różnica
indeksów jest niedodatnia, transwersalne
połączenie różnych punktów jest niemożliwe.
\emph{Warunek Morse'a--Smale'a} żąda właśnie
transwersalności wszystkich takich par.
Powyższy rachunek wyjaśnia rolę warunku;
istnienie odpowiednich perturbacji metryki
oraz zwartość odpowiednich uzupełnień przestrzeni trajektorii
wymagają osobnych argumentów w późniejszej teorii Morse'a.
Sama zerowymiarowość nie dowodzi skończoności liczby
trajektorii. Sam lokalny iloraz nie dowodzi też, że
globalny iloraz przez czas jest rozmaitością Hausdorffa.

\par\smallskip
{\footnotesize\noindent\emph{Dalsza lektura:}
\href{https://www.math.ru.nl/~mueger/diff_notes.pdf}{dowód twierdzenia Sarda},
\href{https://ocw.mit.edu/courses/18-965-geometry-of-manifolds-fall-2004/pages/lecture-notes/}{wykłady MIT o topologii różniczkowej},
\href{https://people.maths.ox.ac.uk/ritter/morse-cambridge/lecture03.pdf}{przecięcia transwersalne},
\href{https://people.maths.ox.ac.uk/ritter/morse-cambridge/lecture04.pdf}{transwersalność parametryczna i funkcje Morse'a}.\par}
